% =============================================================================
% Limitations Section --- Film Production Incentives Gravity Model
%
% Slots into main document under a \section{Limitations} header
% positioned between Results and Conclusion.
% This file begins at \subsection{} level.
%
% Required packages: natbib (for \citep, \citet)
%
% Bibliography keys referenced:
%   borusyak2024, callaway2021, dechaisemartin2020,
%   goodmanbacon2021, sunabraham2021
% =============================================================================

% ----------------------------------------------------------------
\subsection{Data Limitations}
\label{sec:lim_data}
% ----------------------------------------------------------------

Several IMDb measurement issues affect how bilateral pairs are constructed from the underlying film records. First, IMDb's definition of country of origin is opaque: the platform does not document what determines the classification. This introduces measurement error in the importer variable that cannot be fully characterised. Second, for films listing multiple countries of origin, the proportion of production attributable to each country is not recoverable from the data, and budgets are therefore divided equally across origins. Any attempt to construct production weights would introduce considerable additional complexity for limited gain. Third, filming location is recorded as the primary shooting location only, so secondary production locations are systematically undercounted. This third limitation is likely to bias incentive effects downward, since destinations that have specialised in serving as secondary locations will likely be omitted from the data.

Film budgets are missing for approximately 53\% of films in the sample, so the budget regressions are identified on a non-random subsample biased toward larger studio productions for which budget figures are publicly reported. The selection into the budget subsample is plausibly correlated with the production characteristics that also determine incentive responsiveness, so the budget-margin estimates should be interpreted as conditional on the selection mechanism rather than as estimates of the population-wide effect on production spending.

The most significant data limitation arises in comparing the effectiveness of different incentive instrument types. The design of individual incentive schemes is complex, with eligibility criteria, benefits, and requirements that vary substantially across jurisdictions. The Olsberg SPI Global Incentives Index 2025 documents these features but does not lend itself to a clean classification: the high-level three-category typology (rebate, tax credit, tax shelter) used in this analysis necessarily abstracts from the structural design features that determine effective generosity in practice. The Model~4 instrument-type estimates should be interpreted with this caveat in mind. The continuous generosity specification (Model~5) addresses the headline rate but this is a simplication and does not reflect the structural design dimensions.

Subnational incentives are not captured in the data. Many schemes operate at the subnational level. For example, many U.S.\ states administer substantial incentive programmes, but this analysis does not capture them and treats countries like the U.S\ without a national scheme as not having a scheme at all. Notably, this is likely to bias the origin coefficients downward, as domestic producers respond to subnational incentives that are not measured here.

% ----------------------------------------------------------------
\subsection{Identification Limitations}
\label{sec:lim_identification}
% ----------------------------------------------------------------

The country-pair fixed effects design absorbs all time-invariant bilateral heterogeneity but cannot rule out time-varying confounders. This is the central identification limitation of the analysis. If incentive adoption coincides with other policy or structural changes, the pair fixed effects estimator cannot separate the two. For example, the European Union implemented both the MEDIA Incentive Programme and the Audiovisual Media Services Directive, both of which plausibly affected bilateral co-production flows independently of national incentive policy. To the extent that incentive adoption is correlated with such time-varying policy changes, the pair fixed effects estimates conflate the two.

The two-way fixed effects estimator used in Models~8 through~10 may be biased under treatment effect heterogeneity across adoption cohorts. A growing econometric literature \citep{goodmanbacon2021, callaway2021, dechaisemartin2020, sunabraham2021, borusyak2024} has shown that the two-way fixed effects estimator with staggered treatment adoption recovers a weighted average of all possible two-by-two difference-in-differences comparisons in the sample, including comparisons that use already-treated units as controls for later-treated units. When treatment effects vary across cohorts, these comparisons can receive negative weights, biasing the estimated coefficient. The Model~7 finding that early adopters exhibit substantially larger effects than late adopters constitutes direct evidence of cohort-level treatment effect heterogeneity, increasing the likelihood that the two-way fixed effects estimators would produce biased estimates. The Model~9 event study provides information on structure of the treatment response but does not itself resolve the issue, since the relative-time coefficients are themselves weighted averages across cohorts. Re-estimation using cohort-robust estimators such as the imputation approach of \citet{borusyak2024} or the group-time average treatment effect estimator of \citet{callaway2021} would provide a useful robustness check, particularly to test whether the estimated within-pair destination effect of approximately 9\% for film counts is sensitive to the specific weighting structure of the two-way fixed effects estimator.

% ----------------------------------------------------------------
\subsection{Generalisability and External Validity}
\label{sec:lim_generalisability}
% ----------------------------------------------------------------

The sample period is dominated by the pre-streaming theatrical era. The analysis ends in 2023, but the international film and television industry has shifted dramatically over the latter part of the sample window toward streaming originals and global episodic production. Production location decisions in the streaming era are made by a different set of decision-makers operating under different commercial logics than the studio system that dominates the historical sample. 

Estimated late-adopter effects may also overstate the marginal returns available to countries considering adoption today. The Model~7 finding that early adopters attracted substantially more international production than late adopters reflects diminishing marginal returns in an increasingly crowded policy environment. However, a country considering adoption in 2026 faces an even more saturated landscape than the late adopters in this sample did and presumably an even smaller marginal effect of any single additional incentive scheme.
